\section{Finite attainment on a prescribed receiver}
\label{sec:attainment96}
We make the spectral converse operational on the actual initial
Schmidt support.  This also records all normalization factors for a
three-dimensional example.  Neither construction replaces a fixed
pure initial acquisition by an announced mixture.

\begin{proposition}[Support-exact finite instrument]
\label{prop:finiteinstrument96}
Let $C_0:\C^d\to R$ satisfy $C_0^*C_0=\rho$, $\tr\rho=1$, and
let $S_0=\operatorname{ran}C_0\subseteq R$.  Suppose
$\rho=\sum_{h=1}^n w_h a_h$, where $w_h>0$, $\sum_hw_h=1$ and
$\operatorname{rank}a_h\le r$.  For each $h$ let $b_h$ be a density
matrix of rank at most $\ell$.  Choose coefficient maps
\[
 C_h:\C^d\to\C^r,\quad C_h^*C_h=w_ha_h,
 \qquad D_h:\C^d\to\C^\ell,\quad D_h^*D_h=b_h.
\]
The Moore--Penrose inverse $C_0^+:R\to\C^d$, zero on $S_0^\perp$,
gives $V_h=C_hC_0^+:R\to\C^r$ with
\begin{equation}\label{eq:finiteinstrument96}
 V_hC_0=C_h,\qquad \sum_h V_h^*V_h=P_{S_0}.
\end{equation}
These maps extend to an instrument on $R$ with the same $n$ classical
outcomes and quantum outputs of dimension at most $r$.  With the fresh
vector $|D_h\rangle$ prepared independently after outcome $h$, the
receiver block is the product in \eqref{eq:resetblock91}.
For the spectral decomposition used in
Theorem~\ref{thm:spectralvalue94}, $n\le d$ suffices.
\end{proposition}
\begin{proof}
If $x\in\ker\rho$, positivity and the barycenter equation imply
$\langle x,a_hx\rangle=0$ for every $h$.  Hence $C_hx=0$ and
$C_hP_{\operatorname{supp}\rho}=C_h$.  The Moore--Penrose identities
then give $V_hC_0=C_h$ and
\[
 \sum_hV_h^*V_h=(C_0^+)^*\rho C_0^+=P_{S_0}.
\]
This proves \eqref{eq:finiteinstrument96} even when $\rho$ is singular.
To complete the instrument without adding a recorded outcome, choose
a unit vector $e\in\C^r$ and add to its first branch the completely
positive map
\[
             X\longmapsto\tr(P_{S_0^\perp}X)|e\rangle\langle e|.
\]
Together with $X\mapsto V_1XV_1^*$ in that branch and
$X\mapsto V_hXV_h^*$ in the others, the total map is trace
preserving.  This completion vanishes on all physically occupied
initial receiver states.  It is normal also on a larger separable
receiver space, and retains no dilation environment.  The fresh vector
has squared norm $\tr b_h=1$.  Measurement of each input gives the
product block by expansion in the fixed input bases.  The final POVMs
therefore implement the chosen leaf optima.

For the spectral theorem, put $q=q^{(r)}(\lambda(\rho))$.
The least-majorant proof in Lemma~\ref{lem:spectralroof94} gives
$\lambda(\rho)\prec q$.  A decreasing vector majorized by $q$ lies
in the convex hull of its permutations; this follows by successive
two-coordinate averaging, or equivalently by the doubly stochastic
characterization of majorization.  Every permutation belongs to the
affine hyperplane $\sum_i x_i=1$, of dimension $d-1$.  Carath\'eodory
therefore reduces its convex combination to at most $d$ terms.
Use those permutations as diagonal atoms in a fixed eigenbasis of
$\rho$ and apply the construction above.  Eigenvalue ties can change
the chosen eigenbasis or decomposition, not the value or the bound.
The weights may be real; this argument asserts no bit or gate
complexity bound.
\end{proof}
The probability of $(y,h)$ under a first measurement $E^\theta$ is
$w_h\tr[a_h(E_y^\theta)^{\mathsf T}]$.  Thus $w_h$ is generally not
the conditional receiver-outcome probability after observing $y$.
No physical branch is postselected in the attaining construction.

\subsection{One complete three-dimensional calculation}
Let $d=3$, $r=2$, $t=1$, and
\[
 \rho=\diag(3/4,1/8,1/8),\qquad
 |C_0\rangle=\frac{\sqrt3}{2}|1,1\rangle+
            \frac{|2,2\rangle+|3,3\rangle}{2\sqrt2}.
\]
Its squared norm is one.  For a fully specified finite device family,
put $\omega=e^{2\pi i/3}$.  Use the coordinate basis and the three
bases
\[
 |u_j^{(a)}\rangle=\frac1{\sqrt3}
       \sum_{x=0}^2\omega^{a x^2+jx}|x+1\rangle,
       \qquad a=0,1,2,\quad j=0,1,2,
\]
with all six orders of each basis.  These $24$ ordered bases have
weights $1/24$.  Their projections satisfy
\eqref{eq:basismoments90}.  To check the equal-label identity directly,
consider the $12$ vectors before ordering.  For their nine chirp
vectors the $(x,u;x',u')$ tensor entry in the sum of projectors
squared is
\[
 \frac19\sum_{a,j=0}^2
   \omega^{a(x^2+u^2-x'^2-u'^2)+j(x+u-x'-u')}.
\]
It equals one exactly when the two multisets $\{x,u\}$ and
$\{x',u'\}$ agree, and zero otherwise: equal sums and square sums
over $\mathbb F_3$ imply equal products.  The three coordinate
vectors add one to the fully diagonal entries.  The full sum is
therefore $I+S$.  Dividing by $12$ proves the equal-label moment;
summing the other two labels gives $(3I-S)/24$ for distinct labels.
This supplies the exact weighted family, rather than only its existence.

Under the scalar hypothesis $C_y=I/3$.  Under the alternative one of
the $24$ ordered projective measurements is drawn \emph{once} and
used twice.  The priors are $p_C=3/7$, $p_E=4/7$, so each individual
alternative has unconditional prior $1/42$.  All these weights are
rational, whereas the vectors and values below are algebraic.

\begin{proposition}[A prescribed-spectrum message witness]
\label{prop:qutritwitness96}
For this experiment and initial acquisition the exact communicated
and no-message values with old dimension three and fresh dimension
two are
\begin{equation}\label{eq:qutritvalues96}
 P_{3,2}^{\rm b}(\rho)=\frac47+\frac{\sqrt3}{56},\qquad
 P_{3,2}^{\rm nm}(\rho)=\frac47+\frac{\sqrt6}{112}.
\end{equation}
Two receiver outcomes attain the first value using an old output of
dimension two.  The second value retains the complete old receiver
and is an upper over all no-message strategies, not just the exhibited
one.  The strict gain is $(2\sqrt3-\sqrt6)/112$.
\end{proposition}
\begin{proof}
The water-filling rule has $m=1$, $\tau=1/4$, hence
$q^{(2)}=(3/4,1/4,0)$.  Its secular root is $\sqrt3/4$, whereas
the leading two original eigenvalues have root $\sqrt6/8$.
For two positive entries $(a,b)$ the secular equation reduces to
$x^2=ab$.  In particular the disjoint rational enclosures
\[
 \frac{433}{1000}<\frac{\sqrt3}{4}<\frac{434}{1000},\qquad
 \frac{306}{1000}<\frac{\sqrt6}{8}<\frac{307}{1000}
\]
are certified by squaring their positive endpoints.  They are
intervals for algebraic numbers, not exact rational values.
Theorem~\ref{thm:spectralvalue94} and
Corollary~\ref{cor:messagevalue94}, whose no-message reduction is
Lemma~\ref{lem:postponement96}, give both upper bounds in
\eqref{eq:qutritvalues96}.

For explicit attainment use
\[
 a_1=\diag(3/4,1/4,0),\quad
 a_2=\diag(3/4,0,1/4),\quad w_1=w_2=1/2,
\]
so $\rho=(a_1+a_2)/2$.  With $C_0=\sqrt\rho:\C^3\to\C^3$, set
\[
 \begin{array}{ll}
 V_1=\begin{pmatrix}1/\sqrt2&0&0\\0&1&0\end{pmatrix},&
 D_1=\begin{pmatrix}1/\sqrt2&0&0\\0&1/\sqrt2&0\end{pmatrix},\\[6pt]
 V_2=\begin{pmatrix}1/\sqrt2&0&0\\0&0&1\end{pmatrix},&
 D_2=\begin{pmatrix}1/\sqrt2&0&0\\0&0&1/\sqrt2\end{pmatrix}.
 \end{array}
\]
Every displayed map has domain $\C^3$ and codomain $\C^2$.
The old coefficients are $C_h=V_hC_0$, and directly
\[
 \sum_hV_h^*V_h=I_3,\quad C_h^*C_h=a_h/2,\quad
 \tr D_h^*D_h=1.
\]
Prepare $|D_h\rangle$ independently given the receiver outcome.
On $\C^2\otimes\C^2$ let
$|\psi^-\rangle=(|1,2\rangle-|2,1\rangle)/\sqrt2$.
For $y=z$ declare $C$ on
$|\psi^-\rangle\langle\psi^-|$; for $y\ne z$ its effect is zero.
Use the complement for the alternative.  Each branch filtered swap
has negative eigenvalue $-\sqrt3/16$, with eigenvector $|\psi^-\rangle$.
The payoff coefficient is $\kappa=1/21$, so the total improvement,
with three first labels and two outcomes, is
$3\cdot2\cdot\sqrt3/(21\cdot16)=\sqrt3/56$.
The same instrument is used after every first label.

Without a message keep $C_0$ and use $D_1$ after every first label.
The negative eigenvector is the antisymmetric vector in the first two
old coordinates and the two fresh coordinates; its eigenvalue is
$-\sqrt6/16$.  The same equal-label readout gives improvement
$3\cdot\sqrt6/(21\cdot16)=\sqrt6/112$.
The squared-norm normalizations above, not an unnormalized Bell
vector, determine these factors.  All outcomes are included and no
initial mixture or favorable-history selection has been introduced.
\end{proof}
The same construction with $0<t<1$ uses full-rank devices; its gain is
$t^2(2\sqrt3-\sqrt6)/[16(8-t^2)]$.  At $t=0$ it is zero and both
values equal $1/2$.  A positive gain at this fixed $\rho$ does not
contradict the globally optimized no-message attainments at other
initial spectra.  The retained-register interpretation presupposes
the specified calibrated reset cut.
